% Ch. 10 : vérifier, comparer, agir
\documentclass[border=6pt]{standalone}
\usepackage{coursfig}
\begin{document}
\begin{tikzpicture}[x=1mm,y=1mm]
  \node[box=Enc, text width=46mm, minimum height=24mm] (r) at (0,0)
       {\textbf{Lire l'état réel}\sub{le paquet est-il là ?}\sub{le fichier a-t-il ce contenu ?}};
  \node[diamond, aspect=1.7, draw=Ocre, fill=OcreBg, line width=.8pt, align=center, inner sep=1mm] (d) at (62,0)
       {\textbf{Conforme à}\\\textbf{l'état désiré ?}};
  \node[box=Sea, text width=40mm, minimum height=18mm] (ok) at (124,16)
       {\textbf{Ne rien faire}\\[.5mm]{\normalsize rapporter} \cmd{ok}};
  \node[box=Brick, text width=40mm, minimum height=18mm] (ch) at (124,-16)
       {\textbf{Agir, le minimum}\\[.5mm]{\normalsize rapporter} \cmd{changed}};
  \draw[flow] (r) -- (d);
  \draw[flow=Sea] (d.north) |- node[elabel, pos=.25, left=1mm, fill=none, text=Sea]{oui} (ok.west);
  \draw[flow=Brick] (d.south) |- node[elabel, pos=.25, left=1mm, fill=none, text=Brick]{non} (ch.west);
\end{tikzpicture}
\end{document}
